\documentclass[12pt,a4paper]{article}
\usepackage{ugmath}
\usepackage{float}
\usepackage{placeins}
\usepackage{array}
\usepackage[labelfont=bf,labelsep=space]{caption}
\newcommand{\paren}[1]{\left( #1 \right)}
\newcommand{\alumno}{Ricardo León Martínez}
\newcommand{\materia}{Algebra Lineal I}
\newcommand{\profesor}{Claudia Reynoso Alcántara}
\newcommand{\tarea}{Tarea 7}
\newcommand{\fecha}{23/3/2026}

\begin{document}

    \begin{center}
        {\large\textbf{UNIVERSIDAD DE GUANAJUATO}}\\[0.3cm]
        {\normalsize\textbf{DIVISIÓN DE CIENCIAS NATURALES Y EXACTAS}}\\
        {\normalsize\textbf{CAMPUS GUANAJUATO}}\\[1cm]

        {\Large\textbf{\tarea\ (\materia)}}\\[1cm]
    \end{center}

    \noindent
    \textbf{Nombre:} \alumno 
    \hfill 
    \textbf{Fecha:} \fecha 
    \hfill 
    \textbf{Calificación:} \rule{3cm}{0.4pt}

    \vspace{0.3cm}

    En todos los problemas $F$ es un campo que contiene a $\mathbb{Q}$.

    \begin{mdframed}[style=mdbluebox,frametitle={Ejercicio 1}]
        Encuentra una base de $\mathbb{R}^{3}$ que contenga a los vectores $(1,1,1)$ y
        $(1,2,3)$, si existe. Si no existe, demuéstralo.
    \end{mdframed}

    \begin{solucion}
        Sean $a,b\in\mathbb{R}$ tales que
        \begin{equation*}
            a(1,1,1)+b(1,2,3)=(0,0,0).
        \end{equation*}
        Entonces
        \begin{equation*}
            (a+b,a+2b,a+3b)=(0,0,0),
        \end{equation*}
        lo cual implica el sistema
        \begin{equation*}
            \begin{cases}
                a+b=0\\
                a+2b=0\\
                a+3b=0
            \end{cases}
        \end{equation*}
        Restando la primera ecuación de la segunda se obtiene $b=0$,
        y sustituyendo en $a+b=0$ se sigue que $a=0$. Por tanto,
        $\{v_{1},v_{2}\}$ es linealmente independiente. Ahora,
        buscamos un vector $v_{3}=(x,y,z)\in\mathbb{R}^{3}$ tal que $\{v_{1},v_{2},v_{3}\}$
        sea linealmente independiente. Esto ocurre si la única solución de
        \begin{equation*}
            \alpha v_{1}+\beta v_{2}+\gamma v_{3}=(0,0,0)
        \end{equation*}
        es $\alpha=\beta=\gamma=0$. Equivalentemente, $v_{3}$ no debe pertenecer al subespacio
        generado por $v_{1}$ y $v_{2}$, es decir, no debe existir $a,b\in\mathbb{R}$ tales que
        \begin{equation*}
            v_{3}=av_{1}+bv_{2}.
        \end{equation*}
        Tomemos por ejemplo
        \begin{equation*}
            v_{3}=(0,0,1).
        \end{equation*}
        Veamos que $v_{3}\notin\langle\{ v_{1},v_{2}\}\rangle$. Supongamos que existen $a,b\in\mathbb{R}$
        tales que
        \begin{equation*}
            a(1,1,1)+b(1,2,3)=(0,0,1).
        \end{equation*}
        Entonces
        \begin{equation*}
            (a+b,a+2b,a+3b)=(0,0,1),
        \end{equation*}
        lo cual implica
        \begin{equation*}
            \begin{cases}
                a+b=0\\
                a+2b=0\\
                a+3b=1
            \end{cases}
        \end{equation*}
        Restando la primera de la segunda se obtiene $b=0$, de donde $a=0$. Sustituyendo en
        la tercera ecuación se obtiene 0=1, contradicción. Por tanto, $v_{3}\notin\langle\{ v_{1},v_{2}\}\rangle$.
        Así, el conjunto
        \begin{equation*}
            \{(1,1,1),(1,2,3),(0,0,1)\}
        \end{equation*}
        es linealmente independiente en $\mathbb{R}^{3}$, y al tener tres vectores en un espacio
        de dimensión 3, forma una base.
    \end{solucion}

    \begin{mdframed}[style=mdbluebox,frametitle={Ejercicio 2}]
        Encuentra una base del espacio de soluciones en $F^{4}$ del sistema:
        \begin{equation*}
            \begin{cases}
                x_{1}+4x_{2}+3x_{3}-x_{4}=0\\
                3x_{1}-5x_{2}+9x_{3}-2x_{4}=0\\
                2x_{1}-x_{2}+6x_{3}-x_{4}=0.
            \end{cases}
        \end{equation*}
    \end{mdframed}

    \begin{solucion}
        Sea $V\subseteq F^{4}$ el conjunto de soluciones del sistema
        \begin{equation*}
            \begin{cases}
                x_{1}+4x_{2}+3x_{3}-x_{4}=0\\
                3x_{1}-5x_{2}+9x_{3}-2x_{4}=0\\
                2x_{1}-x_{2}+6x_{3}-x_{4}=0.
            \end{cases}
        \end{equation*}
        Su matriz asociada es
        \begin{equation*}
            \begin{pmatrix}
                1 & 4 & 3 & -1\\
                3 & -5 & 9 & -2\\
                2 & -1 & 6 & -1
            \end{pmatrix}
        \end{equation*}
        su forma escalonada y reducida por filas es
        \begin{equation*}
            \begin{pmatrix}
                1 & 0 & 3 & 0\\
                0 & 1 & 0 & 0\\
                0 & 0 & 0 & 1
            \end{pmatrix}
        \end{equation*}
        De aqui tenemos el siguiente sistema de ecuaciones
        \begin{equation*}
            \begin{cases}
                x_{1}+3x_{3}=0\\
                x_{2}=0\\
                x_{4}=0
            \end{cases}
        \end{equation*}
        Así, sea $t\in F$, tomamos $x_{3}=t$. Entonces
        \begin{equation*}
            x_{1}=-3t,\quad x_{2}=0,\quad x_{3}=t,\quad x_{4}=0.
        \end{equation*}
        Por lo tanto, el conjunto solución es de la forma
        \begin{equation*}
            (3t,0,t,0)=t(-3,0,1,0).
        \end{equation*}
        En consecuencia, el espacio de soluciones es
        \begin{equation*}
            V=\langle\{(-3,0,1,0)\}\rangle
        \end{equation*}
        Dado que este vector es no nulo, es linealmente independiente, luego forma una base.
        Así, $\{(-3,0,1,0)\}$ es una base del espacio de soluciones.
    \end{solucion}

    \begin{mdframed}[style=mdbluebox,frametitle={Ejercicio 3}]
        Sean $W_{1},W_{2}$ subespacios vectoriales de dimensión 2 en $F^{3}$, demuestra que
        \begin{equation*}
            \dim_{F}(W_{1}\cap W_{2})\gt0.
        \end{equation*}
    \end{mdframed}

    \begin{proof}
        Supongamos que
        \begin{equation*}
            \dim_{F}(W_{1}\cap W_{2})=0.
        \end{equation*}
        Entonces
        \begin{equation*}
            W_{1}\cap W_{2}=\{0\}.
        \end{equation*}
        Sea $\{u_{1},u_{2}\}$ una base de $W_{1}$ y $\{v_{1},v_{2}\}$ es una base de $W_{2}$.
        Digamos que el conjunto
        \begin{equation*}
            \{u_{1},u_{2},v_{1},v_{2}\}
        \end{equation*}
        es linealmente independiente. En efecto, supongamos que existen $a_{1},a_{2},b_{1},b_{2}\in F$
        tales que
        \begin{equation*}
            a_{1}u_{1}+a_{2}u_{2}+b_{1}v_{1}+b_{2}v_{2}=0.
        \end{equation*}
        Reordenando,
        \begin{equation*}
            a_{1}u_{1}+a_{2}u_{2}=-(b_{1}v_{1}+b_{2}v_{2}).
        \end{equation*}
        El lado izquierdo pertenece a $W_{1}$ y el derecho a $W_{2}$, por lo tanto
        \begin{equation*}
            a_{1}u_{1}+a_{2}u_{2}\in W_{1}\cap W_{2}
        \end{equation*}
        Por la suposición, $W_{1}\cap W_{2}=\{0\}$, asi que
        \begin{equation*}
            a_{1}u_{1}+a_{2}u_{2}=0.
        \end{equation*}
        Como $\{u_{1},u_{2}\}$ es base, se sigue que $a_{1}=a_{2}=0$. Entonces
        \begin{equation*}
            b_{1}v_{1}+b_{2}v_{2}=0
        \end{equation*}
        y como $\{v_{1},v_{2}\}$ es base, resulta $b_{1}=b_{2}=0$. Por lo tanto,
        el conjunto es linealmente independiente. Hemos encontrado 4 vectores linealmente
        independiente en $F^{3}$, lo cual es imposible, ya que cualquier conjunto linealmente
        independiente en $F^{3}$ tiene a lo mas 3 elementos. La suposición es falsa,
        por lo tanto
        \begin{equation*}
            \dim_{F}(W_{1}\cap W_{2})\neq0.
        \end{equation*}
        Dado que la dimensión es un entero no negativo, se concluye
        \begin{equation*}
            \dim_{F}(W_{1}\cap W_{2})\gt0.
        \end{equation*}
    \end{proof}

    \begin{mdframed}[style=mdbluebox,frametitle={Ejercicio 4}]
        Sea $V$ un espacio vecorial de dimensión $n$ sibre un campo $F$. Determina si las siguientes
        afirmaciones son falsa $(\textbf{F})$ o verdaderas $(\textbf{V})$, justifica tus respuestas:
        \begin{enumerate}
            \item[(a)] Para todo $v\in V$, existe una base de $V$ que contiene a $v$.
            \item[(b)] Si $\{v_{1},\dots,v_{n}\}$ genera a $V$, entonces es una base de $V$.
            \item[(c)] Si $S=\{v_{1},\dots,v_{n},v_{n+1}\}$ genera a $V$, entonces existe un subconjunto de $S$ con $n$ elementos
            que es una base de $V$.
            \item[(d)] Si $w\in V\setminus\{0\}$ y $\mathcal{B}=\{v_{1},\dots,v_{n}\}$ es una base de $V$, entonces
            existe $i\in\{1,\dots,n\}$ tal que al reemplazar $w$ por $v_{i}$ en $\mathcal{B}$, seguimos teniendo una base
            de $V$. 
        \end{enumerate}
    \end{mdframed}

    \begin{solucion}
        \begin{enumerate}
            \item[(a)] Para todo $v\in V$, existe una base de $V$ que contiene a $v$.\\
            \textbf{Falso}\\
            Si $v=0$, entonces no puedo pertenecer a ninguna base, pues todo conjunto que contiene
            al vector cero es linealmente dependiente.
            \item[(b)] Si $\{v_{1},\dots,v_{n}\}$ genera a $V$, entonces es una base de $V$.\\
            \textbf{Verdadero}\\
            Sea $\{v_{1},\dots,v_{n}\}$ un conjunto generador de $V$ con $n$ elementos. Supongamos
            que es linealmente dependiente. Entonces existe $i$ tal que $v_{i}$ es combinacion lineal
            de los demás, por lo que
            \begin{equation*}
                \{v_{1},\dots,v_{i-1},v_{i+1},\dots,v_{n}\}
            \end{equation*}
            sigue generando $V$, lo cual contradice que ningun conjunto con menos de $n$ elementos
            puede generar un espacio de dimensión $n$. Por lo tanto es linealmente independiente,
            y asi es una base.
            \item[(c)] Si $S=\{v_{1},\dots,v_{n},v_{n+1}\}$ genera a $V$, entonces existe un subconjunto de $S$ con $n$ elementos
            que es una base de $V$.\\
            \textbf{Verdadero}\\
            El conjunto $S$ tiene $n+1$ elementos, luego es linealmente dependiente. Entonces existe
            algún vector de $S$ que es combinación lineal de los demás. Eliminandolo, obtenemos un
            subconjunto con $n$ elementos que sigue generando $V$. Por el inciso (b), ese subconjunto
            es una base.
            \item[(d)] Si $w\in V\setminus\{0\}$ y $\mathcal{B}=\{v_{1},\dots,v_{n}\}$ es una base de $V$, entonces
            existe $i\in\{1,\dots,n\}$ tal que al reemplazar $w$ por $v_{i}$ en $\mathcal{B}$, seguimos teniendo una base
            de $V$.\\
            \textbf{Verdadero}\\
            Como $\mathcal{B}$ es base, existen únicos escalares $a_{1},\dots,a_{n}\in F$ tales que
            \begin{equation*}
                w=a_{1}v_{1}+\cdots+a_{n}v_{n}.
            \end{equation*}
            Como $w\neq0$, existe al menos un índice $i$ tal que $a_{i}\neq0$. Despejando
            \begin{equation*}
                v_{i}=\frac{1}{a_{i}}\paren{w-\sum_{j\neq i}a_{j}v_{j}},
            \end{equation*}
            por lo que $v_{i}$ pertenece al subespacio generado por
            \begin{equation*}
                \{v_{1},\dots,v_{i-1},w,v_{i+1},\dots,v_{n}\}
            \end{equation*}
            Esto implica que dicho conjunto genera $V$. Además, tiene $n$ elementos, por lo que es
            una base
        \end{enumerate}
    \end{solucion}

    \begin{mdframed}[style=mdbluebox,frametitle={Ejercicio 5}]
        Sea $V$ un espacio vectorial sobre $F$ y sea $\{v_{1},v_{2},v_{3},v_{4}\}$ una base de $V$. Consideramos los
        vectores $w_{1}=v_{1}+v_{2},w_{2}=v_{2}+v_{3},w_{3}=v_{3}+v_{4}$ y $w_{4}=v_{4}+v_{1}$. Encuentra la demensión
        y una base del subespacio de $V$ generado por $\{w_{1},w_{2},w_{3},w_{4}\}$. ¿Existe algún subconjunto de $\{w_{1},w_{2},w_{3},w_{4}\}$ que
        sea una base de este subespacio? justifica tu respuesta
    \end{mdframed}

    \begin{solucion}
        Observemos que
        \begin{equation*}
            w_{1}-w_{2}+w_{3}-w_{4}=(v_{1}+v_{2})-(v_{2}+v_{3})+(v_{3}+v_{4})-(v_{4}+v_{1})=0.
        \end{equation*}
        Por lo tanto $\{w_{1},w_{2},w_{3},w_{4}\}$ es linealmente dependiente, y en consecuencia
        \begin{equation*}
            \dim\langle w_{1},w_{2},w_{3},w_{4}\rangle\leq3.
        \end{equation*}
        Sea $a,b,c\in F$ tales que
        \begin{equation*}
            aw_{1}+bw_{2}+cw_{3}=0.
        \end{equation*}
        Entonces
        \begin{equation*}
            a(v_{1}+v_{2})+b(v_{2}+v_{3})+c(v_{3}+v_{4})=0.
        \end{equation*}
        Agrupando
        \begin{equation*}
            av_{1}+(a+b)v_{2}+(b+c)v_{3}+cv_{4}=0.
        \end{equation*}
        Como $\{v_{1},v_{2},v_{3},v_{4}\}$ es base, se tiene
        \begin{equation*}
            a=0,\quad a+b=0,\quad b+c=0,\quad c=0.
        \end{equation*}
        De $a=0$ se sigue que $b=0$, y de ahí $c=0$. luego
        \begin{equation*}
            a=b=c=0.
        \end{equation*}
        Por tanto, $\{w_{1},w_{2},w_{3}\}$ es linealmente independiente. Como hay
        tres vectores linealmente independientes en el subespacio generado por $\{w_{1},w_{2},w_{3},w_{4}\}$,
        se concluye
        \begin{equation*}
            \dim\langle w_{1},w_{2},w_{3},w_{4}\rangle=3.
        \end{equation*}
        Una base por ejemplo es
        \begin{equation*}
            \{w_{1},w_{2},w_{3}\}
        \end{equation*}
        Veamos que existen subconjuntos que sean base. En efecto, de hecho, cualquier conjunto de tres vectores
        linealmente independiente es base. Ya vimos que
        \begin{equation*}
            \{w_{1},w_{2},w_{3}\}
        \end{equation*}
        es una base. Además, como
        \begin{equation*}
            w_{4}=w_{1}-w_{2}+w_{3},
        \end{equation*}
        cualqueir conjunto de tres de los $w_{i}$ que no contenga redundancia también será base.
    \end{solucion}

    \begin{mdframed}[style=mdbluebox,frametitle={Ejercicio 6}]
        Demuestra que el espacio vectorial $\mathcal{F}=\{f:\mathbb{R}\to\mathbb{R}: f\text{ es una función continua}\}$ no es un
        espacio de dimensión finita sobre $\mathbb{R}$
    \end{mdframed}

    \begin{proof}
        Para cada $k\in\mathbb{N}$, definimos
        \begin{equation*}
            f_{k}(x)=x^{k},\quad x\in\mathbb{R}.
        \end{equation*}
        Cada $f_{k}$ es continua, luego $f_{k}\in\mathcal{F}$. Sea $n\in\mathbb{N}$ y
        consideramos el conjunto
        \begin{equation*}
            \{f_{0},f_{1},\dots,f_{n}\}=\{1,x,x^{2},\dots,x^{n}\}
        \end{equation*}
        Supongamos que existen $a_{0},\dots,a_{n}\in\mathbb{R}$ tales que
        \begin{equation*}
            a_{0}f_{0}+a_{1}f_{1}+\cdots+a_{n}f_{n}=0,
        \end{equation*}
        es decir,
        \begin{equation*}
            a_{0}+a_{1}x+a_{2}x^{2}+\cdots+a_{n}x^{n}=0
        \end{equation*}
        para todo $x\in\mathbb{R}$. Esto significa que el polinomio
        \begin{equation*}
            p(x)=a_{0}+a_{1}x+\cdots+a_{n}x^{n}
        \end{equation*}
        es cero en $\mathbb{R}$. Por tanto, $\{1,x,x^{2},\dots,x^{n}\}$ es linealmente independiente.
        Así, existe un conjunto de $n+1$ funciones linealmente independientes en $\mathcal{F}$
    \end{proof}

    \begin{mdframed}[style=mdbluebox,frametitle={Ejercicio 7}]
        Sea $W=\{(x_{1},\dots,x_{n})\in\mathbb{R}:x_{1}+2x_{2}+\cdots+nx_{n}=0\}$. Demuestra que $W$ un subespacio
        vectorial de $\mathbb{R}^{n}$ y encuentra una base para él. 
    \end{mdframed}

    \begin{proof}
        Veamos que cumple con tres propiedades. Primero contiene al vector cero
        \begin{equation*}
            (0,\dots,0)\in W\quad\text{ pues }\quad0+2\cdot0+\cdots+n\cdot0=0.
        \end{equation*}
        Segundo es cerrado bajo suma. Sea $u=(x_{1},\dots,x_{n})$ y $v=(y_{1},\dots,y_{n})$
        en $W$, entonces
        \begin{equation*}
            x_{1}+2x_{2}+\cdots+nx_{n}=0,\quad y_{1}+2y_{2}+\cdots+ny_{n}=0.
        \end{equation*}
        Sumando
        \begin{equation*}
            (x_{1}+y_{1})+2(x_{2}+y_{2})+\cdots+n(x_{n+y_{n}})=0
        \end{equation*}
        Luego $u+v\in W$. Tercero es cerrado bajo producto escalar. Sea $\lambda\in\mathbb{R}$
        y $u\in W$
        \begin{equation*}
            \lambda x_{1}+ 2\lambda x_{2}+\cdots+n\lambda x_{n}=\lambda(x_{1}+2x_{2}+\cdots+nx_{n})=0.
        \end{equation*}
        Luego $\lambda u\in W$. Por tanto $W$ es subespacio de $\mathbb{R}^{n}$. Ahora encontremos una base,
        la condición es
        \begin{equation*}
            x_{1}+2x_{2}+\cdots+nx_{n}=0.
        \end{equation*}
        Despejamos $x_{1}$
        \begin{equation*}
            x_{1}=-2x_{2}-3x_{3}-\cdots-nx_{n}.
        \end{equation*}
        Entonces cualquier vector de $W$ es
        \begin{equation*}
            (x_{1},\dots,x_{n})=(-2x_{2}-3x_{3}-\cdots-nx_{n},x_{2},x_{3},\dots,x_{n}).
        \end{equation*}
        Separando
        \begin{equation*}
            =x_{2}(-2,1,0,\dots,0)+x_{3}(-3,0,1,0,\dots,0)+\cdots+x_{n}(-n,0,\dots,0,1).
        \end{equation*}
        Ya vimos que cualquier vector de $W$ se escribe como combinación de esos vectores.
        Ahora supongamos
        \begin{equation*}
            a_{2}(-2,1,0,\dots,0)+\cdots+a_{n}(-n,0,\dots,0,1)=(0,\dots,0).
        \end{equation*}
        Comparamos coordenada a coordenada, luego, todos son cero por lo tanto son independientes.
    \end{proof}
\end{document}